% Draft material for main.tex.
% Purpose: document the validation mode in which the POWHEG hardest emission is
% kept, but the QTilde shower reconstruction is bypassed. This file is not
% included by the paper yet; it is a staging note for use once the corresponding
% plots have been produced and checked.

\subsection*{Draft insertion: fixed-order POWHEG validation without shower reconstruction}

\paragraph{Suggested placement.}
Insert this material after
Fig.~\ref{fig:validation-diff-dijet-pol-balance} and before
\texttt{\textbackslash subsection\{Real-Emission Spin Information in the Shower\}}.
At that point the text has just discussed the fixed-order \POLDIS\ comparison
and the instability of the Breit-frame balance observables in the standard
Herwig event-record path. The new paragraph can then clarify that this
instability is not a failure of the real-emission matrix element itself, but a
consequence of passing the accepted POWHEG emission through the shower
reconstruction machinery before it is analysed.

\paragraph{Draft text.}
The jet-balance observables provide a particularly sensitive diagnostic of how
the accepted POWHEG emission is represented in the event record. At fixed order,
and before any shower recoil or reconstruction is applied, the NLO real-emission
configuration contains two hard coloured partons in the Breit current system.
For the dijet observables considered here this implies a sharply balanced
configuration, with $p_{T,2}^{B}/p_{T,1}^{B}$ concentrated at unity and
$A_{p_T}$ concentrated near zero, up to binning and acceptance effects. In the
standard \Herwig\ POWHEG path, however, the accepted hardest emission is handed
to the QTilde shower through the usual hard-tree construction and inverse
reconstruction. This step is required for the physical NLO+PS event, but it
reshuffles the partonic momenta before the event reaches the fixed-order-like
Rivet analysis. The resulting Breit-frame jet-balance distributions are
therefore not a direct comparison of the raw real-emission matrix element with
\POLDIS.

To isolate this effect we introduce a validation-only mode,
\texttt{FixedOrderNoShower}. In this mode the POWHEG hardest emission is still
generated by the same polarized DIS matrix element and with the same accepted
event weight, but the returned real-emission partons are inserted directly into
the event record and the QTilde hard-tree construction, inverse reconstruction,
shower evolution, intrinsic-$p_T$ reconstruction, hadronization, and decays are
bypassed. The POWHEG matrix-element partons retain the diagnostic role tags used
by the \texttt{MC\_DIS\_BREIT} analysis, so the correlated-weight validation can
be run with
\begin{quote}
\small
\path{MC_DIS_BREIT:JETINPUT=MEPARTONS:RIVETWEIGHTS=YES}.
\end{quote}
This construction is not intended as a physical showered prediction. Rather, it
provides a fixed-order closure test of the realized real-emission kinematics and
of the optional helicity weights used to construct $\sigma$ and
$\Delta\sigma$ on common event kinematics.

The comparison shown in Fig.~\ref{fig:validation-fixedorder-noshower-balance}
uses the same DIS acceptance and dijet cuts as the \POLDIS\ validation, but
with the shower reconstruction disabled. In this setup the
$p_{T,2}^{B}/p_{T,1}^{B}$ and $A_{p_T}$ distributions collapse to the expected
fixed-order structures and the leading and subleading Breit-frame jet
transverse-momentum spectra agree with the \POLDIS\ reference within the
statistical and scale uncertainties. This demonstrates that the real-emission
implementation and the correlated helicity-weight construction reproduce the
fixed-order kinematics once the subsequent shower-interface reconstruction is
removed. The broader balance distributions observed in the standard
shower-reconstructed path should therefore be interpreted as an effect of the
NLO+PS event construction, not as a mismatch in the underlying real-emission
matrix element.

\begin{figure}[!htbp]
  \centering
  % Replace these filenames once the final plots are copied into figures/.
  \includegraphics[width=0.47\linewidth]{figures/pT1_fixedOrderNoShower.pdf}
  \includegraphics[width=0.47\linewidth]{figures/pT2_fixedOrderNoShower.pdf}

  \includegraphics[width=0.47\linewidth]{figures/pT2OverpT1_fixedOrderNoShower.pdf}
  \includegraphics[width=0.47\linewidth]{figures/pTAsym_fixedOrderNoShower.pdf}
  \caption{Validation of the raw POWHEG real-emission kinematics with the
  QTilde shower reconstruction disabled. The comparison uses the
  \texttt{FixedOrderNoShower} POWHEG-emission mode together with
  \texttt{MC\_DIS\_BREIT:JETINPUT=MEPARTONS:RIVETWEIGHTS=YES}. The balance
  observables provide a direct check that the two leading Breit-frame
  matrix-element partons reproduce the fixed-order expectation before the
  standard NLO+PS reconstruction is applied.}
  \label{fig:validation-fixedorder-noshower-balance}
\end{figure}

\paragraph{Appendix/workflow text.}
In Appendix~\ref{app:rivet-workflow}, after the paragraph introducing
\path{MC_DIS_BREIT:JETINPUT=MEPARTONS:RIVETWEIGHTS=YES}, add a short workflow
note:

\begin{quote}
\small
For a validation-only fixed-order closure test, the same analysis can be run
with \path{--fixed-order-powheg-no-shower}. The campaign driver then inserts
\path{set /Herwig/Shower/ShowerHandler:POWHEGEmissionMode FixedOrderNoShower}
in the POWHEG NLO cards, keeps
\path{set /Herwig/Shower/ShowerHandler:HardEmission POWHEG}, and disables
hadronization and decays. The cascade handler remains active only to generate
the POWHEG hardest emission. The returned real-emission partons are analysed
directly by \texttt{MC\_DIS\_BREIT}; this mode is therefore used only for
fixed-order validation and is not a replacement for the showered prediction.
\end{quote}

\paragraph{Cross-reference in the shower-spin section.}
In the later subsection
\texttt{\textbackslash subsection\{Real-Emission Spin Information in the Shower\}},
the text should explicitly separate this fixed-order closure mode from the
physical shower-spin study. A suitable insertion point is after the paragraph
ending ``without obscuring the fixed-order validation of the hard process'':

\begin{quote}
\small
The fixed-order no-shower mode introduced above is not used in the following
spin-correlation study. The observables in this section are designed to test
how the spin-density information of the realized POWHEG emission is propagated
through subsequent shower evolution. They therefore require the standard
shower-reconstructed NLO+PS event construction, with the QTilde shower active.
In the \texttt{FixedOrderNoShower} mode the shower reconstruction and subsequent
branching are deliberately disabled, so this mode can validate the hard
real-emission kinematics but cannot probe the spin effects discussed here.
\end{quote}

\paragraph{Cross-reference in the hadronization discussion.}
The hadronization subsection and Appendix~\ref{app:hadronization-extra} should
make the same distinction. A compact sentence can be added near the beginning
of \texttt{\textbackslash subsubsection\{Impact of Hadronization\}}, after the
first sentence:

\begin{quote}
\small
These comparisons are performed only in the standard showered event mode; the
\texttt{FixedOrderNoShower} validation mode disables both the shower
reconstruction and the non-perturbative event evolution, and is therefore not
part of the hadronization study.
\end{quote}

A similar one-sentence reminder can be used in
Appendix~\ref{app:hadronization-extra} if needed:

\begin{quote}
\small
All configurations in this appendix use the full showered event construction
with hadronization enabled; the fixed-order no-shower validation mode is
excluded by construction.
\end{quote}

\subsection*{Checks to complete before inserting this material}

\begin{itemize}
\item Use the restored common DIS window
  $100~\mathrm{GeV}^2 < Q^2 \le 2500~\mathrm{GeV}^2$ and
  $0.2 \le y \le 0.6$ in Herwig, \POLDIS, and \texttt{MC\_DIS\_BREIT}.
\item Run the correlated-weight fixed-order no-shower campaign with scale
  variations, e.g.
\begin{verbatim}
python3 run_validation_campaign.py full \
  --base-dir ~/Projects/HerwigPol/HwPolNotesNew/DISPOL \
  -t rivetweights_noshower \
  --jobs 192 \
  --shards 200 \
  --seed-base 400000 \
  --posnlo-events 400000000 \
  --negnlo-events 40000000 \
  --scale-variations \
  --poldis skip \
  --poldis-refs-campaign <Q2GT100-POLDIS-CAMPAIGN> \
  --poldis-error-mode scale \
  --rivetweights \
  --setup ALL \
  --fixed-order-powheg-no-shower \
  --force-prepare \
  --yoda-merge-tool /home/apapaefs/Projects/Herwig/Herwig-pol-full-python3-rivet4/bin/yodamerge
\end{verbatim}
\item Confirm the generated POSNLO and NEGNLO cards contain
  \path{GenerateRivetWeights Yes},
  \path{POWHEGEmissionMode FixedOrderNoShower},
  \path{LimitEmissions HardOnly}, and
  \path{MC_DIS_BREIT:JETINPUT=MEPARTONS:RIVETWEIGHTS=YES}.
\item Check the diagnostic histograms:
  \path{METaggedOutgoingCount}, \path{METaggedRole2Count}, and
  \path{METaggedRole3Count}. The selected events should be populated by the
  tagged POWHEG matrix-element partons, with role-2 and role-3 entries present
  in the real-emission sample.
\item Compare \path{pT1}, \path{pT2}, \path{pT2OverpT1}, and \path{pTAsym}, and
  the polarized counterparts \path{DpT1}, \path{DpT2},
  \path{DpT2OverpT1}, and \path{DpTAsym}, against the \POLDIS\ NLO reference.
  The balance observables should be sharply concentrated at the fixed-order
  values, unlike the standard \texttt{ShowerReconstructed} path.
\item When inserting the final paper text, add the explicit exclusions above in
  the shower-spin and hadronization sections: the no-shower mode should be
  presented only as a fixed-order validation device, while the spin-effect and
  hadronization studies require the standard full showered event construction.
\item Keep the wording conditional until the final plots are inspected. If the
  agreement is limited to the balance observables but not to all jet spectra,
  replace the sentence claiming agreement of \path{pT1}/\path{pT2} with a more
  restricted statement.
\end{itemize}
